% SPDX-License-Identifier: Apache-2.0
\section{The Decision That Governs the Rest}
\label{sec:e-elaboration}

Two ways to turn this Rust into hardware. They are not compatible, and
almost everything else follows from the choice.

\begin{figure}[H]
\centering
\begin{tikzpicture}[node distance=6mm]
  % Positions are computed from the boxes rather than from a fixed drop.
  % The two middle boxes have multi-line text, so a fixed `below=16mm of
  % src` put the bottom box on top of them: the build stayed green and the
  % figure was wrong.
  \node[boxf,text width=62mm,align=left] (src) {%
    \centerline{\textbf{Rust source}}\\[2pt]
    units, buses, transactions, tags, config};

  \node[boxf,text width=27mm,align=left,below=10mm of src.south west,
        anchor=north west] (rt) {%
    \centerline{\textbf{Runtime}}\\[2pt]
    the program runs and builds a netlist};

  \node[boxf,text width=27mm,align=left,below=10mm of src.south east,
        anchor=north east] (pm) {%
    \centerline{\textbf{Proc macro}}\\[2pt]
    the macro reads the source and emits hardware};

  % Below whichever of the two hangs lower.
  \node[box,text width=62mm,align=center] (out)
    at ($(src.south)!1!(rt.south |- pm.south) + (0,-12mm)$)
    {\textbf{Netlist, then Verilog or VHDL}};

  \draw[fl] (src.south west) ++(8mm,0) -- (rt.north);
  \draw[fl] (src.south east) ++(-8mm,0) -- (pm.north);
  \draw[fl] (rt.south) -- (rt.south |- out.north);
  \draw[fl] (pm.south) -- (pm.south |- out.north);
\end{tikzpicture}
\caption{The two elaboration models. Only the right-hand one can
reinterpret \code{if} on a signal, because only it sees the source.}
\label{fig:elaboration}
\end{figure}

\textbf{Runtime elaboration.} The program runs, builds a netlist through
side effects, and emits Verilog. Chisel~\cite{bachrach2012} works this way,
and for the structural half of this design it would work here unchanged.
The full expressive power of standard Rust remains available for metaprogramming
and hardware generation without requiring procedural macros. However, native
\code{if} on dynamic signals remains impossible, maintaining the constraints
detailed in Section~\ref{sec:e-walls}, and generator errors surface at execution
time rather than during compilation.

\textbf{Proc macro elaboration.} An attribute on the unit parses the item's
own source and emits hardware from it. \code{if} and \code{match} on
signals can be reinterpreted structurally, which removes
Section~\ref{sec:e-walls} entirely, and mistakes are compile errors. The
cost is real and should be stated plainly: the code looks like Rust and
does not mean Rust, and a reader who assumes Rust semantics inside a unit
will be wrong.

\begin{rulebox}[Recommendation]
Proc macro elaboration. Control flow on a signal is the largest remaining
cost of embedding, and this is the only approach that removes it.
\end{rulebox}

The second thesis' argument applies here unchanged. If a toolchain is
affordable to build, it is worth building the one that buys the most, and a
front end over Rust's syntax tree buys the type system, the module system,
the generics, the package manager and the editor tooling along with it.

\textbf{What was built.} \code{\#[lower]} is the proc macro
elaboration, with one departure from the recommendation as stated.
It reads \code{if} and \code{match} as Rust means them rather than
reinterpreting them, and reads a subset that is also plain Rust,
\code{with!}, \code{case!} and \code{select!} among registers,
wires, memories and channel ports, and refuses the rest with a message
of its own. So the cost above,
code that looks like Rust and does not mean Rust, was not paid: a
unit that lowers is the unit that simulates, one file, and the
runtime~\cite{runtime} and examples~\cite{examples} documents state
what the subset is and check every piece of it against a trace under
\code{nvc}. The RISC-V core~\cite{vreteno} is written in that subset
whole. A function under \code{\#[lower]} is the one addition since:
combinational, a few \code{let}s and a value, plain Rust for the
simulation and inlined at every call in a lowered unit of the same
file. The macro reads the file for it, which
\code{Span::local\_file} names, since a proc macro sees one item at
a time. So a step can be short and its pieces named; the core's
serial port is written so.
